\section{Introduction}
Diabetes affects about 589 million adults worldwide, a number projected to reach 853 million by 2050, and more than four in ten people with diabetes are undiagnosed \cite{idf2025}. About one in five people with diabetes has diabetic retinopathy (DR), and the number of adults with DR is expected to rise from 103 million in 2020 to 160 million in 2045 \cite{teo2021}. On the International Clinical Diabetic Retinopathy (ICDR) scale, the disease runs from no DR (grade~0) through mild, moderate and severe non-proliferative DR (NPDR, grades~1--3) to proliferative DR (PDR, grade~4) \cite{wilkinson2003}. Grade~1 is defined by microaneurysms (MA) alone; it causes no symptoms, but it is the stage at which control of glycaemia and blood pressure is most useful. Annual fundus photography can find it, yet specialist graders are scarce exactly where the number of patients grows fastest \cite{idf2025}.

Deep learning (DL) has changed this picture. Convolutional networks trained on large graded collections reached specialist-level accuracy for referable DR \cite{gulshan2016,ting2017}, autonomous systems passed prospective pivotal trials \cite{abramoff2018,ipp2021}, and a system integrated into the national screening programme of Thailand reached a sensitivity of 91.4\% and a specificity of 95.4\% for vision-threatening DR in real time \cite{ruamviboonsuk2022}. The evidence also shows three weaknesses. First, almost all deployed systems target referable or vision-threatening DR, so mild NPDR is deliberately treated as negative \cite{gulshan2016,ipp2021,ruamviboonsuk2022}. Second, performance varies widely between products and sites \cite{lee2021,wang2026}, the reference standard itself is noisy at the boundary between grades~0 and~1 \cite{krause2018}, and pooled analyses report a positive predictive value of only about one half at realistic prevalence \cite{wewetzer2021}. Third, results obtained on public benchmarks can be inflated by methodological faults such as data leakage between training and test data \cite{kapoor2023} and by dataset bias and evaluation on data from the same source \cite{varoquaux2022}. Reporting guidance for prediction models therefore asks for a clear description of data sources, exclusions, and external validation \cite{collins2024}.

For early detection these weaknesses meet in a specific way. The decisive lesion, the microaneurysm, is a few pixels wide in the native photograph, but most models downsample the image to 224--512 pixels; the extent to which early signs survive this step is rarely quantified. Image-quality control is recommended \cite{dai2021,ruamviboonsuk2022}, but whether rejecting poor images interacts with disease severity is rarely reported. Explanations are usually shown as heatmaps, but they are seldom compared with pixel-level lesion annotations \cite{sayres2019}. Finally, the public datasets on which most models are trained and compared are seldom audited for duplicated images and conflicting labels.

The aim of this study is to develop and externally evaluate a lesion-aware DL system for the detection of any DR (grade~$\geq$1), with explicit attention to grade~1, and to audit the public data on which such systems rely. To this end, four public datasets are first audited for duplicates, label conflicts, acquisition shift, image quality and lesion size; a system with a quality gate, a lesion-query attention head and an imbalance-aware loss is then implemented; and it is compared with a conventional baseline on internal and external data with confidence intervals. The study therefore asks whether such a system detects any DR, and grade~1 in particular, more accurately than a global-pooling CNN baseline, and whether this advantage holds on data from other populations and cameras. It further asks whether attention maps that are built into the model point to annotated lesions more often than post-hoc Grad-CAM maps, and what the quality gate and the imbalance-aware loss contribute, given the limit that input resolution places on the visibility of early lesions.

Compared with the cited work, the study (i) reports a data audit that found duplicated images with conflicting grades in two public test sets, including one pair that crosses the official IDRiD train/test split, and measures how small microaneurysms become at typical input resolutions; (ii) shows that image-quality rejection is strongly associated with severe disease, a safety issue that is not discussed in screening studies; (iii) evaluates heatmaps quantitatively against pixel masks with a chance baseline; and (iv) releases the code and split definitions. The work was run on a CPU-only budget, so it is a reduced-scale implementation of the protocol whose negative results are reported as such.

The rest of the article is organised as follows. Section~\ref{sec:lr} reviews the literature and appraises each key study; Section~\ref{sec:mm} describes data, system and statistics; Section~\ref{sec:res} reports the data audit and the experimental results; Section~\ref{sec:conc} concludes, followed by the data and code availability statements.

\section{Literature Review}\label{sec:lr}
The review is organised by the role each study plays for this work; Table~\ref{tab:lr} gives the strengths, weaknesses and contribution of every key study. Four studies older than five years are retained deliberately: Gulshan et al.\ \cite{gulshan2016} and Ting et al.\ \cite{ting2017} are the original validations against which all later systems are compared, Abr\`amoff et al.\ \cite{abramoff2018} is the first prospective trial that led to regulatory authorisation, and Sayres et al.\ \cite{sayres2019} is still the only randomised reader study of heatmap assistance for DR. No newer work replaces these roles; all other review sources date from 2021 or later, and no conference papers are used in this section.

\textbf{Screening systems and their limits.} The landmark studies \cite{gulshan2016,ting2017} established expert-level detection of referable DR on large internal and multi-ethnic external data, and the pivotal trial of Abr\`amoff et al.\ \cite{abramoff2018} moved autonomous screening into primary care. All three, however, optimise the referable threshold and return a single image-level label. Lee et al.\ \cite{lee2021} then showed that seven commercial algorithms differed widely in sensitivity on the same 311{,}604 images, and the meta-analyses of Wewetzer et al.\ \cite{wewetzer2021} and Wang et al.\ \cite{wang2026} confirm high pooled accuracy but a low positive predictive value and few fully external validations. For this study they define the target operating point (sensitivity and specificity of about 0.90) and the requirement of external testing.

\textbf{Lesion-aware and quality-aware models.} DeepDR \cite{dai2021} combines quality assessment, lesion segmentation and grading and detects mild NPDR, but it needs lesion masks for about 15{,}000 images. CABNet \cite{he2021} adds category and global attention blocks so that every grade receives its own discriminative feature channels, which counters class imbalance without lesion masks, and CRA-Net \cite{ma2024} adds a transformer-guided lesion-relation attention evaluated on DDR and APTOS. These journal studies motivate the weakly supervised, attention-based head used here; their limitation is that attention maps were judged visually rather than against masks. The foundation models RETFound \cite{zhou2023} and DeepDR Plus \cite{dai2024} show that self-supervised pretraining on hundreds of thousands of fundus photographs improves label efficiency and allows progression prediction, but they require compute and data far beyond a student project.

\textbf{Surveys and explainability.} Recent surveys \cite{tsiknakis2021,bhulakshmi2024} agree that public datasets carry only image-level grades, that classes are strongly imbalanced, and that early disease and lesion localisation remain open problems; multimodal models are expected to help when clinical data are available \cite{qin2026}. The reader study of Sayres et al.\ \cite{sayres2019} showed that heatmaps can even reduce accuracy for patients without DR, which is why localisation must be measured, not assumed.

\textbf{Gap.} No reviewed study (a) targets grade~1 explicitly while auditing the data used for evaluation, (b) quantifies how many microaneurysms remain visible at the input resolution, or (c) tests attention maps against lesion masks with a chance baseline. This study addresses these gaps.

\begin{table}[t]
\centering\scriptsize
\caption{Critical appraisal of the key studies (ordered as in the text).}\label{tab:lr}
\begin{tabular}{@{}p{1.9cm}p{3.4cm}p{3.4cm}p{3.4cm}p{2.9cm}@{}}
\toprule
Study & Design and data & Strengths & Weaknesses & Contribution here\\\midrule
Gulshan 2016 \cite{gulshan2016} & Inception-v3 ensemble, $>$128k images; EyePACS-1, Messidor-2 & First expert-level referable-DR result; two validation sets & Referable DR only; image-level output; internal-like data & Baseline concept and performance reference\\
Ting 2017 \cite{ting2017} & 76k images, 10 multi-ethnic external cohorts & Broad external validation & Referable DR; deployment untested & Requirement of external testing\\
Abr\`amoff 2018 \cite{abramoff2018} & Prospective trial, 900 patients, 10 sites & Pre-specified endpoints; regulatory approval & More-than-mild DR; one product & Target sensitivity and specificity\\
Lee 2021 \cite{lee2021} & 7 algorithms, 311k images, veterans & Head-to-head real-world test & No lesion-level analysis & Expect performance drop externally\\
Dai 2021 \cite{dai2021} & DeepDR: quality + lesion + grading, $\sim$0.5M images & Detects mild NPDR; real-time quality control & Needs $\sim$15k lesion masks & Quality gate design\\
He 2021 \cite{he2021} & CABNet attention blocks on 3 datasets & Handles imbalance; few parameters & Maps judged visually & Attention-based head\\
Ma 2024 \cite{ma2024} & CRA-Net, DDR and APTOS & Models lesion relations & No external test; no mask evaluation & Lesion-relation attention\\
Zhou 2023 \cite{zhou2023} & RETFound, 904k CFP, self-supervised & Label-efficient, generalises & Very high compute; EyePACS in pretraining & Future backbone; overlap check\\
Dai 2024 \cite{dai2024} & DeepDR Plus, progression in 5 years & Personalised screening intervals & Chinese data only & Motivation for early detection\\
Sayres 2019 \cite{sayres2019} & Randomised reader study, 10 graders & Measures effect of explanations & One heatmap method & Heatmaps must be validated\\
Wewetzer 2021 \cite{wewetzer2021}; Wang 2026 \cite{wang2026} & Meta-analyses, primary care / approved devices & Pooled accuracy; sources of heterogeneity & Heterogeneous reference standards & PPV at $\pi=0.10$; external data\\
Tsiknakis 2021 \cite{tsiknakis2021}; Bhulakshmi 2024 \cite{bhulakshmi2024} & Surveys of DL for DR & Datasets, metrics, open problems & Mostly descriptive & Imbalance-aware loss; per-grade metrics\\
\bottomrule
\end{tabular}
\end{table}

\section{Materials and Methods}\label{sec:mm}
\subsection{Design and deviations from the protocol}
The study is a retrospective diagnostic-accuracy study reported along TRIPOD+AI \cite{collins2024}: a system is developed on one source and evaluated against a reference standard on a held-out internal set and on external sets from other countries. The primary task is any DR (grade~$\geq$1 vs.\ grade~0); five-class grading is secondary. Compared with the planned protocol, six deviations were necessary and are stated here. (1) The RETFound weights could not be obtained in the computing environment and a ViT-Large cannot be fine-tuned on the available CPU, so the proposed head was implemented on an ImageNet EfficientNet-B0 \cite{tan2019} and the foundation-model baseline was dropped. (2) Messidor-2 and the graded part of DDR could not be downloaded; DDR is represented by its 757 lesion-annotated images. (3) APTOS~2019 has no patient identifiers, so a patient-level split is impossible; it was moved from the development pool to external testing, where the split problem does not arise. (4) Development used a random subset of 6{,}000 EyePACS patients. (5) The optional clinical branch was removed because no dataset contains clinical variables. (6) Because of the computing time available, each condition was trained once (seed~0, 5 epochs) instead of with several seeds, and the plain cross-entropy (A2) and 448-pixel (A3) ablations were not run; uncertainty is expressed through bootstrap confidence intervals, and the effect of resolution is analysed through the lesion-size audit.

\textbf{Overlap with foundation-model pretraining.} RETFound was pretrained on 904{,}170 colour fundus photographs, 90.2\% from the Moorfields MEH-MIDAS dataset and 9.8\% from Kaggle EyePACS \cite{zhou2023}; APTOS, IDRiD and Messidor-2 were only its downstream evaluation sets. The external sets of this study are therefore not part of that pretraining corpus, whereas EyePACS is; any future RETFound-based version must not report EyePACS internal results as unseen data.

\subsection{Data}
Five public, de-identified sources were used (Table~\ref{tab:data}). \emph{EyePACS} \cite{cuadros2009} (USA; Kaggle 2015 training release, 35{,}126 graded images of 17{,}563 patients) provided the development data: 6{,}000 patients were drawn at random (seed~0), and their 12{,}000 images were split by patient 70/15/15 into training, validation and internal test sets. \emph{EyeQ} \cite{fu2019} provides good/usable/reject quality labels for 12{,}543 EyePACS training images; only labels of training-split patients were used to train the quality gate. \emph{APTOS~2019} \cite{aptos2019} (India, several cameras, 3{,}662 images) and \emph{IDRiD} \cite{porwal2018} (India, Kowa VX-10$\alpha$, 516 images, official train and test parts pooled) are external test sets with ICDR grades. \emph{DDR} \cite{li2019} (China) contributes 757 images with pixel masks of microaneurysms (MA), haemorrhages (HE), hard exudates (EX) and soft exudates (SE); all contain lesions, so DDR measures sensitivity and lesion localisation.

\textbf{Preprocessing and audit.} Every image was cropped to its circular field of view (FOV; pixels with maximal channel value $>$20), padded to a square, resized to 512 pixels, and described by its mean brightness, RMS contrast and sharpness (variance of the Laplacian) inside the FOV. Near-duplicates were searched with a 255-bit perceptual hash of the high-pass filtered green channel, which encodes the vessel pattern. The distribution of nearest-neighbour Hamming distances was clearly bimodal (all duplicate pairs at $\leq$10 bits, the next pair at 13 bits and the bulk at 20--90 bits), so pairs at $\leq$10 bits were treated as duplicates; visual inspection confirmed them. In external sets, one copy of a duplicate group was kept when all grades agreed and the whole group was excluded when they conflicted. In EyePACS, duplicates and both eyes of a patient were assigned to the same split. MA, HE, EX and SE masks of DDR were split into connected components, and each component's equivalent diameter was rescaled to the 224- and 448-pixel model inputs.

\subsection{Proposed system and baseline}
\textbf{Quality gate.} An EfficientNet-B0 with the partial fine-tuning described below and a three-class head was trained on EyeQ labels (3{,}004 images). Images classified as \emph{reject} are not graded but returned for re-capture or referral.

\textbf{Backbone.} The image is resized to $224\times224$ pixels and normalised with ImageNet statistics. Stages 1--6 of an ImageNet-pretrained EfficientNet-B0 \cite{tan2019} are frozen and their output ($192\times7\times7$) is cached; stage~7 and the $1\times1$ convolutional head (1{,}280 channels) are fine-tuned. This choice keeps training feasible on two CPU cores; random horizontal and vertical flips of the cached feature maps serve as augmentation.

\textbf{Lesion-aware head.} The $N=49$ feature vectors $X\in\mathbb{R}^{N\times d}$ ($d=1280$) are layer-normalised, and $Q=8$ learnable lesion queries $P\in\mathbb{R}^{Q\times d_k}$ ($d_k=256$) attend to them:
\begin{equation}
A=\mathrm{softmax}\!\left(\frac{P\,(XW_K)^{\top}}{\sqrt{d_k}}\right)\in\mathbb{R}^{Q\times N},\qquad F=A\,XW_V .
\end{equation}
Each row of $A$ is a spatial map; the maximum over queries is the lesion heatmap, which is produced in the forward pass and is therefore intrinsic. The query features are averaged and passed to a five-grade head and an any-DR head. The loss is $L=L_{\mathrm{focal}}+L_{\mathrm{BCE}}+0.1\,L_{\mathrm{div}}$, where the focal loss \cite{lin2020} of the grade head uses $\gamma=2$ and class weights inversely proportional to class frequency, $L_{\mathrm{BCE}}$ is the binary cross-entropy of the any-DR head, and $L_{\mathrm{div}}$ is the mean pairwise cosine similarity of the query maps, which keeps the queries apart.

\textbf{Conditions.} \emph{B1} (baseline) uses the same trainable backbone with global average pooling, dropout and a five-class cross-entropy head; its any-DR score is $1-p(\text{grade }0)$, and its heatmaps are Grad-CAM maps \cite{selvaraju2020} of this score. \emph{P} is the proposed system with the quality gate. \emph{A1} is P without the gate (all images graded).

\textbf{Training.} AdamW, one-cycle schedule with peak learning rates $3\times10^{-4}$ (backbone) and $10^{-3}$ (head), weight decay 0.05, batch 64, 5 epochs, seed 0. The checkpoint with the best validation quadratic weighted kappa is kept. The any-DR threshold is the highest threshold reaching a validation sensitivity of 0.90; it is fixed before any test set is used.

\subsection{Evaluation and statistics}
Sensitivity (Se) and specificity (Sp) for any DR at the fixed threshold, the area under the ROC curve (AUC), the positive predictive value at a prevalence of $\pi=0.10$, $\mathrm{PPV}_{0.10}=\mathrm{Se}\,\pi/(\mathrm{Se}\,\pi+(1-\mathrm{Sp})(1-\pi))$, the quadratic weighted kappa $\kappa_w$, the macro-F1 score and grade-1 recall (share of grade-1 images flagged as DR) were computed. 95\% confidence intervals (CI) come from 1{,}000 bootstrap resamples, drawn by patient in EyePACS and by image elsewhere \cite{efron1993}. Paired AUCs were compared with the DeLong test \cite{delong1988}, and sensitivities and specificities with the exact McNemar test; for all other differences the paired bootstrap CI is reported. Localisation is measured with the pointing game \cite{zhang2018}: a hit is scored when the maximum of the upsampled heatmap lies within one 16-pixel patch (at 224 pixels) of an annotated lesion. Because large lesions cover much of some images, the expected hit rate of a random point inside the FOV is reported as a chance baseline. Success criteria fixed in the protocol were: external Se $\geq0.90$ with Sp close to 0.90; grade-1 recall of P significantly above B1; $\kappa_w\geq0.80$; and a pointing-game score of P above Grad-CAM.
